\documentclass[11pt]{article}

\usepackage{fullpage}

\begin{document}

\title{ARM Checkpoint... }
\author{Mahdi Begg\and Arafath Hamid\and Umar Tauqir\and Zakariyah Bello}

\maketitle

\section{Group Organisation}

We split Part I based on the main parts of the emulator and the different instruction types:

\begin{itemize}
    \item Branch instructions
    \item Data processing instructions
    \item Load/store instructions
    \item Registers and memory
    \item Pipeline and instruction decoding
    \item I/O
    \item Testing
\end{itemize}

For communication, we mainly used GitLab, Discord, and WhatsApp. Everyone worked on their own branch so we didn’t interfere with each other’s code. Once each module was working and tested, we merged it into the develop branch. This helped reduce merge conflicts and made it easier to keep track of progress and allowed us to work on the emulator in parallel.

We kept the repository organised by creating separate files for different parts of the emulator before everyone started developing in their own branches. This made it easier to understand which part of the code each person was responsible for. When one part depended on another, such as branch execution needing access to the PC or decoded instruction fields, we communicated this so that the interfaces would match.

\section{Group Progress}

The group is working well overall. Since everyone has a clear task, we can work at the same time without disrupting another person's work via conflicts on the repository. We communicate heavily when there are dependencies from one part of the code to another, which helps avoid confusion.

One thing that has worked well is that each person can test the part they made before it is merged. This makes it easier to find where a bug is coming from, because the code is still isolated before it becomes part of the full emulator. After merging, we can then run the provided tests and our own smaller unit tests to check that the different modules still work together.

For later tasks, the group may need to coordinate more closely before starting implementation. The assembler will probably need more shared planning because parsing assembly text, handling labels, and encoding instructions all depend on a consistent structure. If each person designs their part differently, merging could become harder than it was for the emulator. To avoid this, we will decide as a group on the main data/file structures to lay out before writing code.

\section{Emulator Structure}

We structured the emulator in a modular way, with separate components for memory, registers, bit utilities, pipeline/decode structures, and instruction categories such as branch, data processing, and load/store. This allowed group members to work on separate parts while keeping the overall design organised.

Some parts should be reusable for the assembler, especially the bit utility functions and instruction format knowledge. The assembler will need to perform the reverse process of the emulator: instead of decoding binary instructions, it will encode assembly instructions into binary. Therefore, utilities for bit manipulation should be useful later.

The emulator is structured around the general execution pipeline of fetching, decoding, and executing instructions. The pipeline and decoding code is responsible for identifying the instruction type, while the separate instruction modules handle lower level decoding and execution involving machine state update. This separation helps keep the code readable because the main emulator does not need to contain every instruction case directly.

For the assembler, we expect the bit utilities and instruction formatting to be reusable. The execution functions will not be reused directly, because the assembler is not executing instructions. However, the same understanding of fields such as sf, opc, rn, rd, rm, immediate values, branch offsets, and load/store addressing modes will be needed when encoding assembly instructions into binary.

\section{Implementation Challenges}

One challenge we expect is combining everything together properly. Since different people are working on different parts, making sure everything fits together (like decoding, execution, and updating registers) could be tricky. We are trying to deal with this by merging often and testing after integrating changes.

Another challenge is making sure all instructions are implemented correctly, especially edge cases like branching or memory addressing. To deal with this, we are using both the provided tests and our own tests.

The assembler will likely be harder, since it involves parsing text and converting it into the correct binary format. This is more complex than the emulator because it involves handling labels, offsets, and syntax. To avoid problems later, we plan to spend more time planning the structure before starting to implement it.

For the assembler, labels and offsets are likely to be one of the main challenges. The assembler will need to know where each label is and calculate the correct offset and then encode it into the correct number of bits. We will soon encounter the assembler and plan it effectively so as to reduce the need to redesign later in the development phase.

\end{document}